% Pursuing Stacks (cotes 134-2 to 134-8): archive cote 134-3, site page 25 (PDF page 26), typescript page 25.
% typescript. Draft compared: ps2.tex:1160-1224.
% AI-transcribed by Claude Opus 5.5 (claude-opus-5-5, adaptive thinking, effort high), Message Batches API, batch ps1 (msgbatch_01Rv5hCaNTtUt4JHx1yqoHPL), 2026-11-12; not yet checked by eye.
\scanpage{134-3/25}
\xout{\unsure{25}}
\mspage{25}

this class.  Maybe even it could be shown that this equivalence class
can be recovered in terms of the corresponding functor (*),  in the same
way as (according to Giraud) a ``topology'' on M can be recovered from
the associated subtopos of $M^{\wedge}$\hand{,}namely the associated category of sheaves
on M (in such a way that the set of \xout{\unsure{subcategor}}``topologies''on M can be
identified with the set of \xout{\unsure{subt}}``closed subtopoi'' of $M^{\wedge}$).  The very best
one could possibly hope for along these lines would be a one-to-one
correspondance between isomorphism classes of functors (*)(satisfying
certain properties ?),  and the so-called ``Hot-equivalence classes'' of
topologies on M.

Whether or not this tentative hope is excessive, when it comes to
the (still somewhat vague) question of understanding ``which algebraic
structure species are suitable for expressing homotopy types'',  it
might not be excessive though to expect that in all \inserted{such} cases,  M should
be equipped with just \uline{one} ``natural'' Hot-equivalence class of topologies
\margin{\illegible{1}}on M\deleted{\illegible{1}}\hand{,}which moreover (one hopes,  or wonders) should be expressible
directly in terms of the \inserted{intrinsic} structure of M,  or,  what amounts to the same,
in terms of the \inserted{full} subcategory $\underline{B}$ of objects of ``finite presentation'',  \xout{\unsure{givi}}
giving rise to M via the equivalence
\[ M \simeq \mathrm{Ind}(\underline{B}) . \]
As we just saw in the case  M = (Cat),  the so-called ``canonical topolo-%
gy'' on M need not be within the natural Hot-equivalence class  -  and I
am at a loss for the moment to give a plausible intrinsic caracterizati-%
on of the latter,  in terms of the category M = (Cat). \hand{(${}^{22'}$).}

\over{\unsure{25}}{26}.  What comes to mind though is that the categories such as \hand{$\Delta$} , \hand{$\square$}
and their analogons (corresponding to multicomplexes rather than mono-%
complexes, for instance\over{\unsure{,}}{)}  can be viewed as (generally not full) subcate-%
gories of M  (in fact,  even of the smaller category (Ord) of all orde\unsure{\hand{red}}
sets).  The topologies we found on M were in fact associated in an evi-%
dent way to the choice of such subcategories.  As was already felt by
the end of the reflection two days ago (p.21),  these subcategories
(denoted there by $\underline{B}_{oo}$  ) have rather special features  -  they are
associated to \inserted{simultaneous} cellular decompositions of spheres of all dimensions  -
and it is this feature,  presumably,  that makes the associated ``trivial''
algebraic structure\inserted{species}, giving rise to the category of set-theoretic
realizations  $\underline{B}_{oo}{}^{\wedge}$ ,  eligible for ``describing homotopy types''.  In the
typical example M = (Cat) though,  contrarily to what was suggested \hand{by}
\margin{mainly}\deleted{on page 21} \inserted{the end of section 23} (when thinking \inserted{mainly} of the rather special although important cases
when M is expressible  as a category  $B^{\wedge}$\hand{,} for some category B such as
